% chapters/theory/01.tex

\chapter{From Text to Transformers}
\label{chap:from-text-to-transformers}

Modern language AI begins with a deceptively simple question: how can a machine process language in a way that preserves meaning, order, and context? The Transformer is the architecture that answered this question at scale. To understand why it works—and why it displaced earlier sequence models—we will follow a sequence of ideas that build on one another: what NLP tasks look like, how text becomes tokens, how tokens become vectors, why recurrence struggled, and how self-attention resolves the core bottlenecks.

\section{Natural Language Processing as a set of problems}
Natural Language Processing (NLP) studies how to analyze and generate human language with algorithms. A practical way to organize the field is to group problems into three families.

In \emph{classification}, text is the input and a label is the output. Sentiment analysis is the canonical example: given a review, predict whether it is positive or negative. Other applications include intent detection (“create an alarm”), language identification, and topic classification. These tasks are commonly evaluated with accuracy, but accuracy can be misleading when classes are imbalanced. In such cases, precision and recall become essential: precision measures how many predicted positives are correct, while recall measures how many true positives are captured. The F1 score summarizes the trade-off between precision and recall via their harmonic mean.

In \emph{sequence labeling}, the model assigns multiple labels, often at the token or span level. Named Entity Recognition (NER) labels spans as \texttt{PERSON}, \texttt{LOCATION}, \texttt{TIME}, and so on. Part-of-speech tagging and syntactic parsing are closely related tasks: they reflect the structure of language rather than a single global label.

In \emph{generation}, both input and output are text. Translation maps a source sentence to a target sentence; question answering maps a question to an answer; summarization maps a document to a shorter document; code generation maps a specification to executable code. Generation is difficult to evaluate because many outputs may be valid. Overlap-based metrics such as BLEU and ROUGE compare a prediction to one or more references, but they can penalize paraphrases that are semantically correct yet lexically different. Perplexity is another common metric: it measures how surprised a language model is by text (lower is better), but it does not directly measure usefulness or correctness for downstream tasks.

\section{Tokenization: turning language into discrete units}
Neural networks operate on numbers, not on raw strings. The first bridge from text to computation is \emph{tokenization}, the process of dividing text into discrete units called tokens. Tokens are the objects that will be embedded, processed, and predicted.

A straightforward approach is word tokenization. It is simple, but it creates immediate challenges: vocabulary size becomes large; unseen words appear at inference time; and morphological variants are treated as unrelated symbols. The words \texttt{bear} and \texttt{bears}, for instance, are nearly identical in meaning, yet word-level tokenization assigns them unrelated indices.

A more robust approach is \emph{subword tokenization}. Subword methods split words into reusable pieces so that common roots and affixes are shared across tokens. This reduces out-of-vocabulary risk and allows the model to reuse learned structure across variants. The trade-off is that sequences become longer, and sequence length is a critical driver of computational cost in attention-based models.

At the extreme, one can tokenize by characters. Character tokenization avoids out-of-vocabulary problems and handles misspellings naturally, but sequences become very long and semantics must be built from very small units. In practice, subword tokenization is the most common compromise: it balances robustness and efficiency.

\section{From tokens to vectors: why one-hot is not enough}
Once text has been tokenized, each token must be represented numerically. A purely symbolic representation is one-hot encoding: if the vocabulary size is $V$, a token is represented by a vector in $\mathbb{R}^V$ with a single 1 at its index and zeros elsewhere. One-hot encoding is technically valid, but it is not semantically informative. All one-hot vectors are orthogonal, so similarity measures cannot express that some words are related and others are not.

The goal is instead to learn \emph{embeddings}: dense vectors whose geometry reflects meaning. A common similarity measure is cosine similarity,
\[
\cos(u,v)=\frac{u\cdot v}{\|u\|\|v\|},
\]
which emphasizes the angle between vectors rather than their magnitudes. Cosine similarity is not a perfect notion of semantics, but it is a useful proxy: good embedding spaces place related tokens closer together.

\section{Learning embeddings via proxy objectives}
How can embeddings be learned from data? A powerful idea is to train a model on a \emph{proxy task} that forces it to internalize linguistic structure. Word2vec is a canonical example. It trains on context prediction: either predict the center word from surrounding words (CBOW) or predict surrounding words from the center word (skip-gram). The proxy task is not the final goal; the learned internal representation is.

A simple neural-network view makes this concrete. Suppose the input is a one-hot vector over vocabulary size $V$. The model maps it through a hidden layer of size $d \ll V$ and then projects back to a $V$-dimensional output, producing a distribution over tokens via softmax. Training uses cross-entropy loss and backpropagation. After training, the hidden representation (or the corresponding weight matrix) serves as the embedding. Intuitively, if the model can predict context well, it must encode meaningful relationships among tokens.

Two practical details matter. First, tokenization interacts with inference-time vocabulary: unseen words in word-level tokenization are commonly mapped to a special unknown token (\texttt{UNK}), collapsing distinct unseen words into a single representation. Subword tokenization reduces this problem substantially. Second, generation typically stops when the model emits an end-of-sequence token (\texttt{EOS}), a convention used across many language modeling systems.

\section{Sequence models: RNNs and long-range dependencies}
Token embeddings alone do not capture word order or contextual meaning within a sentence. A naive fix is to average token embeddings, but averaging discards order. A more principled approach is to use sequence models.

Recurrent Neural Networks (RNNs) process tokens sequentially and maintain a hidden state summarizing what has been seen so far. Conceptually,
\[
h_t = f(h_{t-1}, x_t),
\]
where $x_t$ is the token embedding and $h_t$ is the hidden state at time $t$. This introduces order sensitivity and provides contextual representations for each position.

RNNs can be used for sentence classification by taking the final hidden state, and for token labeling by using the per-token hidden states. They can also be used for generation by predicting tokens step by step. However, classic RNNs struggle with long sequences. Information from the distant past can become inaccessible due to vanishing (or exploding) gradients in backpropagation through time. Even when memory is improved (as in LSTMs, which introduce a cell state designed to retain information), RNN computation remains inherently sequential. This limits parallelism and makes training slow on long contexts.

These two limitations—weak long-range dependency handling and poor parallelization—created strong pressure for a different mechanism.

\section{Attention and self-attention}
Attention provides direct access to relevant tokens without relying on a single sequential hidden state. The key idea is to compute, for each position, a weighted combination of other positions based on relevance. This becomes especially powerful when applied within a single sequence: \emph{self-attention}.

Self-attention is commonly described through three vectors. A \textbf{query} encodes what a token is looking for, a \textbf{key} encodes what a token offers for matching, and a \textbf{value} encodes the content to aggregate. We compare the query to all keys to obtain weights, then apply these weights to values.

In matrix form, scaled dot-product attention is:
\[
\mathrm{Attention}(Q,K,V)=\mathrm{softmax}\!\left(\frac{QK^\top}{\sqrt{d_k}}\right)V.
\]
The $\sqrt{d_k}$ scaling stabilizes dot products as the dimensionality grows. Crucially, this computation is expressed in matrix multiplications, making it highly parallelizable on GPUs.

\section{The Transformer architecture}
The Transformer is a stack of attention and feedforward layers designed to replace recurrence. In its original form it has an encoder and a decoder.

The encoder processes the input sequence using self-attention and feedforward sublayers, producing contextual representations for each input token. The decoder generates output tokens autoregressively. It uses \emph{masked self-attention} so each position attends only to previously generated tokens, and it uses \emph{cross-attention} to attend to the encoder outputs when solving sequence-to-sequence tasks such as translation.

Self-attention does not encode order by itself, so Transformers inject positional information into token representations. Early designs add positional encodings to token embeddings; modern variants introduce relative or rotary schemes, discussed later in the book. What matters here is the conceptual necessity: without a positional signal, self-attention cannot distinguish permutations of the same tokens.

\section{Multi-head attention}
Multi-head attention repeats attention multiple times in parallel, each time with different learned projections. Each head learns its own linear maps to queries, keys, and values; the head outputs are concatenated and projected back to the model dimension. The conceptual motivation is expressivity: different heads can attend to different relationships among tokens. In practice, nothing forces heads to specialize, but the optimization process often produces diverse patterns.

\section{Label smoothing}
Next-token prediction is inherently ambiguous: multiple continuations can be plausible. Label smoothing reduces overconfidence by softening the one-hot training target. Instead of assigning probability 1 to the observed next token and 0 to all others, the correct token receives probability $1-\epsilon$ and the remaining mass $\epsilon$ is distributed across other tokens. This can improve robustness in tasks like translation, where many continuations may be valid.

\section{End-to-end view: from tokens to generation}
Putting the pieces together, a text input is tokenized and mapped to token embeddings. Positional information is injected so the model can represent order. Within each Transformer block, embeddings are projected to $Q$, $K$, and $V$, attention weights are computed via scaled dot-products and softmax, and the result is combined with feedforward transformations. In generation settings, decoding proceeds token by token: the model produces a distribution over the vocabulary at each step, a decoding strategy selects the next token, and generation stops at \texttt{EOS}.

This architecture solved two foundational problems simultaneously: it created direct long-range interactions via attention, and it made training highly parallel through matrix computations. These properties explain why the Transformer became the backbone of modern language models and why most later innovations build on its core components rather than replacing them entirely.

\section{Chapter summary}
Tokenization converts raw text into discrete units suitable for modeling. Embeddings convert these discrete units into vectors where similarity can reflect meaning, and proxy tasks such as context prediction provide a scalable way to learn such representations. Sequence models such as RNNs and LSTMs introduced order sensitivity but struggled with long-range dependencies and parallelization. Self-attention provides direct token-to-token interaction and can be implemented efficiently with matrix multiplications. The Transformer architecture organizes self-attention, feedforward layers, masking, and positional information into a practical system that supports both understanding and generation at scale.

\section{Exercises}
\paragraph{1. Tokenization trade-offs.}
Give an example where word-level tokenization produces an \texttt{UNK} token at inference time. Explain how subword tokenization can reduce this failure mode, and what trade-off it introduces.

\paragraph{2. One-hot similarity.}
Prove that two distinct one-hot vectors are orthogonal. What does this imply about cosine similarity and semantic similarity?

\paragraph{3. RNN limitations.}
Explain in your own words why RNN training is hard to parallelize across time. How does self-attention change this?

\paragraph{4. Attention computation.}
Given $Q \in \mathbb{R}^{n\times d_k}$ and $K \in \mathbb{R}^{n\times d_k}$, what is the shape of $QK^\top$? Interpret its entries.

\paragraph{5. Masked self-attention.}
Why must a decoder use masking during generation? Describe what would go wrong if a decoder could attend to future tokens during training.
